\documentclass[11pt,letterpaper]{article}
\usepackage[T1]{fontenc}
\usepackage[utf8]{inputenc}
\usepackage{lmodern}
\usepackage[margin=1in,headheight=15pt]{geometry}
\usepackage{amsmath,amssymb,amsthm,mathtools}
\usepackage{microtype,booktabs,longtable,array,enumitem}
\usepackage{xcolor,fancyhdr,xurl}
\usepackage[colorlinks=true,linkcolor=blue!45!black,citecolor=green!35!black,urlcolor=blue!45!black]{hyperref}
\usepackage[nameinlink,noabbrev]{cleveref}
\hypersetup{pdftitle={The Set-Theoretic Content of Bounded Surreal Arithmetic},pdfauthor={Research article prepared with ChatGPT},pdfsubject={Birthday cutoffs, hereditary set universes, bi-interpretation, axiom spectra and extension rigidity}}
\pagestyle{fancy}\fancyhf{}
\fancyhead[L]{\small Bounded surreal arithmetic and set theory}
\fancyhead[R]{\small\thepage}
\renewcommand{\headrulewidth}{0.3pt}
\setlength{\parskip}{0.25em}
\setlength{\emergencystretch}{3em}
\setlist{nosep,leftmargin=2em}
\clubpenalty=10000\widowpenalty=10000\displaywidowpenalty=10000
\allowdisplaybreaks[2]\numberwithin{equation}{section}
\newtheorem{theorem}{Theorem}[section]
\newtheorem{lemma}[theorem]{Lemma}
\newtheorem{proposition}[theorem]{Proposition}
\newtheorem{corollary}[theorem]{Corollary}
\theoremstyle{definition}
\newtheorem{definition}[theorem]{Definition}
\newtheorem{example}[theorem]{Example}
\theoremstyle{remark}
\newtheorem{remark}[theorem]{Remark}
\newcommand{\No}{\mathbf{No}}
\newcommand{\Ord}{\operatorname{Ord}}
\newcommand{\bd}{\operatorname{b}}
\newcommand{\TC}{\operatorname{TC}}
\newcommand{\B}{\mathcal B}
\newcommand{\K}{\mathcal K}
\newcommand{\C}{\mathcal C}
\newcommand{\I}{\mathcal I}
\newcommand{\J}{\mathcal J}
\newcommand{\Code}{\operatorname{Code}}
\newcommand{\Bit}{\operatorname{Bit}}
\newcommand{\Pre}{\operatorname{Pre}}
\newcommand{\Vis}{\operatorname{VisOrd}}
\newcommand{\Coll}{\operatorname{Coll}}
\newcommand{\SG}{\operatorname{SignGraph}}
\newcommand{\Etern}{\mathsf E}
\newcommand{\Pack}{\mathsf P}
\newcommand{\Pow}{\mathcal P}
\newcommand{\om}{\omega}
\newcommand{\nsum}{\mathbin{\oplus}}
\newcommand{\nprod}{\mathbin{\otimes}}
\newcommand{\card}[1]{\lvert#1\rvert}
\newcommand{\dom}{\operatorname{dom}}
\newcommand{\ran}{\operatorname{ran}}
\newcommand{\rank}{\operatorname{rank}}
\newcommand{\ceilcard}[1]{\kappa(#1)}
\newcolumntype{P}[1]{>{\raggedright\arraybackslash}p{#1}}
\begin{document}
\begin{titlepage}
\centering
\vspace*{0.5in}
{\large\scshape A foundational research article\par}
\vspace{0.35in}
{\Huge\bfseries The Set-Theoretic Content\par of Bounded Surreal Arithmetic\par}
\vspace{0.28in}
{\Large Hereditary-size universes, coherent bi-interpretations,\par axiom spectra, and extension rigidity\par}
\vspace{0.35in}
{\large September 22, 2026\par}
\vspace{0.12in}
{Research article prepared with ChatGPT\par}
\vspace{0.35in}
\begin{minipage}{0.92\textwidth}
\begin{abstract}
Equip a surreal field with its birthday function. Its global connection with
set theory is already the subject of a bi-interpretation theorem announced by
Chen--Hamkins--Yang. This article investigates a more quantitative question:
\emph{which set universe is recovered from a particular birthday cutoff?}

For an epsilon number $\lambda$, let $\B_\lambda$ be the ordered field
$\No_{<\lambda}$ with birthday, and let $\kappa(\lambda)$ be the least cardinal
not below $\lambda$. We give an explicit interpretation whose collapse is
exactly $H_{\kappa(\lambda)}$. If $\lambda$ is a cardinal, this is a uniform
bi-interpretation with $H_\lambda$. Otherwise it is a uniform
bi-interpretation with $(H_{\kappa(\lambda)},\in,\lambda)$: the missing cutoff
is itself recoverable as an interpreted set. Thus even
$\No_{<\varepsilon_0}$, with birthday, recovers all hereditarily countable sets.
Small certificates for arithmetic justify the reverse interpretation without
assuming Replacement inside the hereditary-size universe.

The resulting elementary-inclusion classification is sharp: for epsilon
numbers $\lambda<\eta$, the inclusion $\B_\lambda\prec\B_\eta$ holds exactly
when both cutoffs are cardinals and $H_\lambda\prec H_\eta$. We also calibrate
two explicit growth principles: birthday eternity corresponds to Replacement
at cardinal cutoffs, while packing birthday fibers corresponds to the strong
limit property. Their conjunction selects precisely the worldly cardinal
cutoffs. Finally we give a power-set-sized detector for new subsets in outer
models, prove extension rigidity at strong-limit cutoffs, and explain the
orientation obstruction to bi-interpretation for surcomplex numbers.
\end{abstract}
\end{minipage}
\vfill
\begin{minipage}{0.92\textwidth}\small
\textbf{Status.} Written proofs of proposed local refinements, not a refereed
publication or a proof-assistant-certified development. The global
bi-interpretation theorem, the classical field-cutoff theorem, and ordinary
well-founded coding are prior work, not discoveries claimed here. Novelty is
limited to the refinements identified in the provenance discussion; an
exhaustive priority claim is not made.
\end{minipage}
\end{titlepage}
\pagenumbering{roman}
\tableofcontents
\clearpage\pagenumbering{arabic}

\section{The question, the results, and their provenance}
\label{sec:intro}

The field structure of the surreal numbers and their construction history
carry very different information. Real-closed-field quantifier elimination
makes all inclusions between real closed subfields elementary in the ordered
field language. A birthday predicate remembers much more: it can distinguish
individual sign coordinates and thereby encode sets and well-founded
relations. The issue addressed here is not merely whether this coding is
possible globally, but how its strength changes when only birthdays below one
ordinal are allowed.

\subsection{What is already known}

Chen, Hamkins, and Yang announced that the surreal ordered field with birthday
order is bi-interpretable with the set-theoretic universe, and described a
first-order theory of surreal arithmetic intended to be bi-interpretable with
ZFC. Hamkins's September and November 2025 slides give the global coding by
pointed well-founded extensional relations, together with birthday-growth
principles; the March 2026 seminar announcement repeats the global theorem
\cite{CHYKobe,CHYNotreDame,CHYCUNY}. This article does \emph{not} claim that
theorem, its basic coding idea, or the association of growth with set-existence
axioms as new. We use a different, size-explicit implementation to study
cutoffs. We do not rely on the full axiomatics of the announced theory being
finalized.

The classical cutoff result is also essential prior work: if $\lambda$ is an
epsilon number, $\No_{<\lambda}$ is a real closed subfield. The original
van den Dries--Ehrlich paper must be used together with its erratum
\cite{vdDE,vdDEerr}; Bournez--Guilmant state the corrected field and
real-closedness results explicitly in their Theorems 1.1--1.2
\cite{BG}. The sign construction and recursive arithmetic used below are
classical \cite{Conway,Gonshor}. Rangel--Mariano provide a separate algebraic
set-theoretic program and foundational discussion; their different language
and aims must not be conflated with birthday-enriched field structures
\cite{RM}.

\subsection{The local questions answered here}

We use $\TC(\{x\})$ for the transitive closure containing $x$ and all its
iterated members, and write
\[
 H_\kappa=\{x:\card{\TC(\{x\})}<\kappa\}.
\]
The following statements are the principal refinements developed in this
article. The word ``uniform'' means that fixed finite formulas work at every
cutoff in the stated class, not that a separate language is chosen for each
cardinal.

\begin{theorem}[Cardinal rounding and cutoff recovery; proved below]
\label{thm:main-rounding}
Let $\lambda$ be an epsilon number and let
\[
 \ceilcard{\lambda}=\min\{\kappa:\kappa\text{ is a cardinal and }\lambda\le\kappa\}.
\]
There is one parameter-free interpretation $\I$ such that
\[
 \I(\B_\lambda)\cong(H_{\ceilcard{\lambda}},\in).
\]
If $\lambda$ is a cardinal, $\B_\lambda$ and $(H_\lambda,\in)$ are uniformly
bi-interpretable. If $\lambda$ is not a cardinal, $\B_\lambda$ and
$(H_{\ceilcard{\lambda}},\in,\lambda)$ are uniformly bi-interpretable.
In the latter case the named ordinal $\lambda$ is definable, as an interpreted
set, in $\B_\lambda$ itself.
\end{theorem}

This is a cardinal-rounding statement, not a rank-cutoff statement.
For example, $\kappa(\varepsilon_0)=\omega_1$, whereas
$V_{\varepsilon_0}$ is not the recovered universe. At a singular cardinal
$\kappa$ the recovered structure is still $H_\kappa$; no regularity assumption
is hidden in the theorem.

\begin{theorem}[Elementary-inclusion spectrum; proved below]
\label{thm:main-elementary}
For epsilon numbers $\lambda<\eta$,
\[
 \B_\lambda\prec\B_\eta
 \quad\Longleftrightarrow\quad
 \lambda,\eta\text{ are cardinals and }H_\lambda\prec H_\eta.
\]
The field reducts, in contrast, always form an elementary inclusion.
\end{theorem}

Two further results explain the set-existence and extension content of this
classification. We introduce fully specified principles $\Etern$ and $\Pack$
in \cref{sec:growth}, rather than silently choosing among possible versions of
the announced surreal-arithmetic axioms.

\begin{theorem}[Growth spectrum; proved below]
\label{thm:main-growth}
For an epsilon-number cutoff, birthday eternity $\Etern$ fails at every
noncardinal cutoff. At an uncountable cardinal $\kappa$,
\[
 \B_\kappa\models\Etern
 \quad\Longleftrightarrow\quad
 H_\kappa\models\text{Replacement}.
\]
Fiber packing $\Pack$ holds exactly at the uncountable strong-limit cardinal
cutoffs. Consequently
\[
 \B_\lambda\models\Etern+\Pack
 \quad\Longleftrightarrow\quad
 \lambda\text{ is a cardinal and }V_\lambda\models\mathrm{ZFC}.
\]
\end{theorem}

\begin{theorem}[Extension rigidity at a strong-limit cutoff; proved below]
\label{thm:main-extension}
Let $M\subseteq N$ be transitive models of ZFC with the same ordinals, and let
$\kappa$ be an uncountable cardinal in both, strong limit in $M$. Then
\[
 \B_\kappa^M\prec\B_\kappa^N
 \quad\Longleftrightarrow\quad
 \B_\kappa^M=\B_\kappa^N
 \quad\Longleftrightarrow\quad
 \Pow^M(\alpha)=\Pow^N(\alpha)\quad(\alpha<\kappa).
\]
Without the strong-limit hypothesis, a new subset of $\alpha<\kappa$ still
prevents elementarity whenever $(2^{\card\alpha})^M<\kappa$.
\end{theorem}

The proofs supply explicit witnesses and specify the sense in which the last
cardinal bound is sharp. It is sharp for the old-power-set \emph{parameter},
not a claim that every other possible detector needs that much space.

\subsection{Repository comparison and novelty boundary}

The repository inspection was pinned to
\begin{center}\small\texttt{4f2645599121fa872c7104995e47f86f0382351f}.\end{center}
The report catalog and the foundations report were inspected, with particular
attention to sign presentations, local recursion, universes, birthday
cutoffs, and size distinctions. The newer surcomplex-automorphism report's
scope was also checked \cite{RepoCatalog,RepoFoundations,RepoAutomorphisms}.
These sources already cover classical cutoff closure, foundational size
obstructions, localization of arithmetic, and automorphism rigidity. None of
those broad subjects is novel merely because it reappears below.

The proposed contribution is the \emph{specific package} of cardinal rounding,
recovery of a noncardinal cutoff as a set, the exact elementary-inclusion
classification, the singular-cardinal-safe growth calibration, and the
explicit extension detector. These formulations were not located in the
inspected sources. That is a bounded literature-and-repository observation,
not proof of worldwide absence or an assertion about unpublished work of
other researchers. In particular, the refinements may overlap with details
of the Chen--Hamkins--Yang project not present in the public announcements.
No named longstanding conjecture is advertised as resolved.

\section{Foundations and the two ordinal languages}
\label{sec:foundations}

\subsection{Ambient theory and class conventions}

The proofs take place in ZFC. All cutoff structures are sets. The optional
statements about the full surreal class are definable-class statements,
understood one first-order formula at a time. Neither a set of all surreals
nor an internally available truth predicate for the entire universe is
assumed. Quotient interpretations use definable equivalence relations; the
Mostowski collapse supplies set representatives when we compare them with
$H_\kappa$. Global Choice is not used.

When a transitive model $M$ is explicitly mentioned, its surreals are the sign
sequences belonging to $M$, and cardinalities with superscript $M$ are computed
there. A statement proved internally in an arbitrary model of ZFC does not
assert that the model's well-founded relations are externally well founded.
Our outer-model comparison theorems explicitly use transitive models.

\subsection{The carrier and the birthday function}

A surreal is a function $s:\alpha\longrightarrow\{-,+\}$ for an ordinal
$\alpha$. Comparison is lexicographic, with termination between $-$ and $+$.
The all-plus sequence of length $\alpha$ is the embedded ordinal $\alpha$.
Throughout, $\bd(s)$ is the \emph{surreal ordinal} corresponding to the length
of $s$. Thus it is a function on the surreal carrier, not an extra ordinal
sort. Set-theoretic lengths and their surreal embeddings are distinguished
by context.

For an epsilon number $\lambda$, define
\[
 K_\lambda=\No_{<\lambda},\qquad
 \B_\lambda=(K_\lambda;0,1,+,-,\cdot,<,\bd).
\]
The language includes no exponentiation, support predicate, membership
predicate, or second-order quantifier.

\begin{lemma}[Birthday order and birthday function]
\label{lem:birthday-language}
On each $\B_\lambda$, the birthday function and the birthday-precedence
relation $x\prec_b y\iff\bd(x)<\bd(y)$ are definitionally interconvertible.
Moreover, the embedded ordinals are exactly the fixed points of $\bd$.
\end{lemma}
\begin{proof}
Given the function, precedence is immediate. Conversely, equality of birthdays
is expressed by neither $x\prec_b y$ nor $y\prec_b x$. In the class of sign
sequences of length $\alpha$, the largest in numerical order is the all-plus
sequence. Therefore $\bd(x)$ is the unique largest member of $x$'s
birthday-equivalence class. A sign sequence is fixed by this function exactly
when it is all plus.
\end{proof}

We abbreviate $\Ord(a)$ by $\bd(a)=a$. Quantification over embedded ordinals
in $K_\lambda$ is precisely quantification over $\alpha<\lambda$.

\subsection{Natural arithmetic, not ordinary ordinal arithmetic}

For ordinal arguments, surreal addition and multiplication are the
Hessenberg operations
\[
 \alpha+\beta=\alpha\nsum\beta,\qquad
 \alpha\beta=\alpha\nprod\beta.
\]
They are strictly increasing in each positive argument where appropriate,
and $\alpha+1$ agrees with ordinary ordinal successor. We use ordinary
ordinal exponentiation only in metamathematical bounds and in explicitly
set-theoretic constructions. It is not silently identified with a field
operation.

\begin{lemma}[Closure and cardinal estimates]
\label{lem:card-estimates}
Every uncountable cardinal is an epsilon number. For infinite ordinals
$\alpha,\beta$, natural sums and products have cardinality at most
$\max(\aleph_0,\card\alpha,\card\beta)$. In particular, if $\kappa$ is an
uncountable cardinal and $\alpha,\beta<\kappa$, their natural sum and product
are below $\kappa$.
\end{lemma}
\begin{proof}
In Cantor normal form, a natural sum joins finitely many exponent terms; a
natural product uses finitely many natural sums of their exponents. The
cardinality estimate follows from the corresponding estimate for
$\omega^\gamma$: for infinite $\gamma$, its cardinality is at most
$\max(\aleph_0,\card\gamma)$. One can see the latter by representing ordinals
below $\omega^\gamma$ by finite Cantor-normal-form lists with exponents below
$\gamma$ and natural-number coefficients. These observations also cover
finite arguments after adjoining $\aleph_0$ to the bound.

If $\kappa$ is uncountable and initial, $\omega^\alpha<\kappa$ for every
$\alpha<\kappa$. Continuity gives $\omega^\kappa\le\kappa$. The opposite
inequality follows from $\omega^\alpha\ge\alpha$ and taking suprema. Hence
$\omega^\kappa=\kappa$.
\end{proof}

For a general epsilon number, closure of natural arithmetic below $\lambda$
also follows from the classical field-cutoff theorem. This simple closure is
all the ordinal arithmetic needed by our relation codes.

\subsection{Hereditary size rather than rank}

\begin{lemma}[Basic hereditary-size facts]
\label{lem:hereditary}
Let $\kappa$ be an uncountable cardinal.
A sign sequence of length $\alpha$ belongs to $H_\kappa$ exactly when
$\alpha<\kappa$. Every subset of an ordinal $\alpha<\kappa$ belongs to
$H_\kappa$. Also $H_\kappa\subseteq V_\kappa$.
\end{lemma}
\begin{proof}
The transitive closure of the usual set-coded graph of a sign sequence has
cardinality at most $\max(\aleph_0,\card\alpha)$, and it contains enough of
the domain ordinal to have cardinality at least $\card\alpha$ when $\alpha$
is infinite. The same estimate applies to a subset of $\alpha$.

If $\card{\TC(\{x\})}\le\mu$ for an infinite $\mu<\kappa$, the membership
relation on this transitive closure is a well-founded relation of size at
most $\mu$. Its rank function takes values below $\mu^+$. This follows by
well-founded induction, using regularity of the successor cardinal $\mu^+$:
the supremum of at most $\mu$ ordinals below $\mu^+$ remains below $\mu^+$.
Thus $\rank(x)<\mu^+\le\kappa$, and $x\in V_\kappa$.
\end{proof}

The statement that every subset of $\alpha$ belongs to $H_\kappa$ does
\emph{not} say that the collection $\Pow(\alpha)$ is an element of
$H_\kappa$. This distinction drives the growth results.

\section{Small certificates for surreal arithmetic}
\label{sec:certificates}

A reverse interpretation cannot simply say ``perform the usual set-theoretic
recursion inside $H_\kappa$.'' At a singular cutoff, Replacement may fail in
that structure. We instead prove that each individual arithmetic computation
has a certificate of the same hereditary cardinal scale as its inputs.

\subsection{Legal cuts can be checked using prefixes}

For sets of sign sequences $L<R$, let $\{L\mid R\}$ denote their simplest
separator. The classical simplicity theorem supplies this separator and the
bound
\begin{equation}
 \bd(\{L\mid R\})\le
 \sup\{\bd(a)+1:a\in L\cup R\}.
 \label{eq:cut-bound}
\end{equation}
The supremum of the empty set is zero. These statements can equivalently be
proved by the sign-by-sign cut construction \cite{Gonshor}.

\begin{lemma}[Prefix test for a cut value]
\label{lem:cut-test}
A sign sequence $z$ equals $\{L\mid R\}$ if and only if it separates $L$ and
$R$ and no proper sign prefix of $z$ separates them.
\end{lemma}
\begin{proof}
The forward implication is immediate. For the converse, suppose $y$ is a
separator shorter than $z$. If $y$ is a prefix of $z$, the stated test fails.
If it is not, their longest common prefix lies numerically between them and
is a proper prefix of $z$. The set of separators is convex, so this common
prefix is another separator. Again the test fails. Thus $z$ has minimum
birthday among separators, and uniqueness of the simplest separator gives
the result.
\end{proof}

The check in this lemma quantifies only over the displayed option sets and
the prefixes of $z$. In set coding it is a bounded check: one does not
quantify over all sign sequences of earlier birthdays.

\subsection{Well-founded option graphs of small size}

A labeled option graph has a set $D$ of vertices, left and right successor
relations, and a rank label that strictly decreases along every option edge.
A valuation of such a graph assigns a sign sequence to each vertex so that
its value passes \cref{lem:cut-test} for the values at its options. When the
graph is numeric, this valuation exists and is unique.

\begin{lemma}[Small numeric graphs have small certificates]
\label{lem:graph-size}
Suppose $\mu$ is infinite and a numeric, well-founded option graph has at
most $\mu$ vertices, with its structural data of hereditary size at most
$\mu$. Then it has a rank-labeled valuation certificate of hereditary size
at most $\mu$.
\end{lemma}
\begin{proof}
The graph rank of every vertex is below $\mu^+$, by the successor-cardinal
argument in \cref{lem:hereditary}. Applying \eqref{eq:cut-bound} by induction
shows that a vertex value has birthday at most that vertex's graph rank.
Thus each rank label and each sign value has hereditary size at most $\mu$.
There are at most $\mu$ vertices and at most $\mu^2=\mu$ edges, so the union
of all these data still has hereditary size at most $\mu$. The existence
construction here takes place in the ambient ZFC universe, not in a
possibly deficient $H_\kappa$.
\end{proof}

\subsection{The computation graph for addition and multiplication}

Take all sign prefixes of the inputs, including the inputs themselves. With
\[
 \mu=\max(\aleph_0,\card{\bd(x)},\card{\bd(y)}),
\]
this prefix family and its data have hereditary size at most $\mu$.
For addition, use finite signed words of prefix atoms $A(a)$; an option moves
one atom to one of its sign-prefix options, reversing left and right for a
negative atom. The root word is $A(x)+A(y)$. This is the usual disjunctive-sum
construction of the numeric game for $x+y$.

For multiplication, use finite signed words of pair atoms $P(a,b)$, with
$a,b$ prefixes of the respective inputs. An option of a positive pair atom
is a signed three-term word
\begin{equation}
 P(a',b)+P(a,b')-P(a',b').
 \label{eq:pair-option}
\end{equation}
It is a left option when $(a',b')$ is left--left or right--right, and a right
option for the two mixed choices. Negative atoms exchange the sides and
negate the replacement word. An option of a word replaces one atom by one
of these words. The root is $P(x,y)$. This spells out Conway's multiplication
recursion without first constructing a large universe of games.

There are at most $\mu$ possible finite words. These graphs are well founded.
For addition, use the natural sum, over atoms of the word, of
$\omega^{\bd(a)}$. For multiplication put
\[
 \rho(a,b)=\bd(a)\nsum\bd(b)
\]
and use the natural sum of $\omega^{\rho(a,b)}$ over the pair atoms.
Every pair atom in \eqref{eq:pair-option} has strictly smaller $\rho$ than the
replaced atom. A finite sum of lower powers of $\omega$ is smaller than the
replaced power, so each move strictly decreases the displayed measure.
The standard numeric-game arithmetic theorem identifies the values of these
graphs with $x+y$ and $xy$; their options are exactly the classical recursive
options, including the signs in the negative summands \cite{Conway}.

\begin{theorem}[Hereditary-size arithmetic certificates]
\label{thm:certificates}
There are fixed first-order set-theoretic formulas defining surreal addition,
negation, and multiplication on sign sequences such that, for inputs of
hereditary size at most an infinite $\mu$, a correct computation has a
certificate of hereditary size at most $\mu$. Verification of a supplied
certificate is absolute between transitive sets containing its data.
Consequently, for every uncountable cardinal $\kappa$, these formulas define
in $H_\kappa$ exactly the external surreal operations on its sign sequences.
No assumption $H_\kappa\models\mathrm{Replacement}$ is required.
\end{theorem}
\begin{proof}
For $+$ and multiplication, existentially quantify a set of word vertices
containing the root and closed under the explicit option rule, its exact
option relations, an ordinal rank labeling decreasing on edges, and a
valuation by sign sequences passing the prefix cut test at every vertex.
The node labels specify the input prefixes and the finite signed words; all
checks of option generation quantify over those input prefix families and
over entries of the finite words. Exactness requires both that each generated
option is present and that every recorded edge is generated by the rule.
The predicates for sign sequence, prefix, numerical comparison, finite word,
and decreasing ordinal rank have bounded verifications on the supplied
sets. In particular, a supplied prefix family is not accepted merely because
its members are prefixes: for each ordinal $\xi\le\dom(s)$ it must contain
the unique restriction $s\mathbin{\upharpoonright}\xi$, and every recorded
prefix must be one of these restrictions. The same completeness check is
imposed on the prefix families of vertex values used by the cut test.
These are bounded checks against $\dom(s)+1$ and the supplied families.
A word length is a finite ordinal; this is checked by requiring every
ordinal at most that length to be zero or a successor. Thus a missing
prefix or a nonstandard-length word cannot evade an option requirement.
The verifier is consequently a fixed bounded set-theoretic predicate,
after including these domains and prefix families as certificate data.
Their cardinality and hereditary-size bounds remain $\mu$.

Any valid certificate is sound: external well-founded induction on its
rank labels identifies each vertex value with the value of its generated
numeric game. This also proves uniqueness of the root value. A certificate
exists by the preceding graph construction and \cref{lem:graph-size}, with
hereditary size at most $\mu$. Negation simply reverses all signs and has the
same size estimate.

For inputs in $H_\kappa$, choose the infinite bound $\mu<\kappa$ given above.
The external certificate belongs to $H_\kappa$; conversely any certificate
there is externally sound by bounded absoluteness. Hence the existential
formulas have the same intended values. The uses of cardinal regularity were
at $\mu^+$, not at $\kappa$.
\end{proof}

\begin{remark}[Scope of the certificate theorem]
This is not a proof that every class-recursive surreal operation is absolute
in every hereditary-size structure. It proves the specific finite-signature
claim needed here. Inversion is definable from multiplication by $xy=1$;
its existence in the epsilon-number cutoffs, and real closedness, are supplied
by the classical cutoff theorem. No normal-form birthday estimates or
unproved sharp product-birthday conjecture enter the argument.
\end{remark}

\section{Definable sign access and a bounded pairing map}
\label{sec:bits}

\subsection{Recovering the sign tree from the two orders}

For $x,y,z$ let $\operatorname{Between}(z;x,y)$ say that $z$ lies in the closed
numerical interval with endpoints $x$ and $y$. Define
\begin{equation}
 \Pre(x,y)\quad\Longleftrightarrow\quad
 \forall z\bigl((\operatorname{Between}(z;x,y)\land z\ne x)
                 \longrightarrow\bd(x)<\bd(z)\bigr).
 \label{eq:prefix}
\end{equation}

\begin{lemma}[Definable prefix relation]
\label{lem:prefix}
In every $\B_\lambda$, \eqref{eq:prefix} holds exactly when $x$ is a sign
prefix of $y$, allowing equality.
\end{lemma}
\begin{proof}
If $x$ is a prefix of $y$, every other point between them must continue the
sign sequence of $x$; a point diverging earlier would lie outside that
interval. Its birthday is therefore larger.

If $x$ is not a prefix of $y$ and $y$ is a proper prefix of $x$, the endpoint
$z=y$ violates the formula. In the remaining case the longest common prefix
of $x,y$ lies strictly between them and has birthday smaller than that of
$x$. This prefix belongs to the same cutoff. Finally, for $x=y$ the formula
is vacuous, as required.
\end{proof}

Define the bit predicate by
\begin{align}
 \Bit(s,\alpha)\quad\Longleftrightarrow\quad &
 \Ord(\alpha)\land\alpha<\bd(s)\ \land\nonumber\\
 &\exists t\bigl(\Pre(t,s)\land\bd(t)=\alpha\land t<s\bigr).
 \label{eq:bit}
\end{align}
The unique prefix of length $\alpha$ lies below $s$ precisely when the sign
of $s$ at position $\alpha$ is plus. Thus
\begin{equation}
 \{\alpha<\delta:\Bit(s,\alpha)\},\qquad\bd(s)=\delta,
 \label{eq:fiber-subsets}
\end{equation}
runs through \emph{all} subsets of $\delta$, once each. This uses a
first-order quantifier over numbers of one birthday, not a second-order
quantifier in the language.

\subsection{Pairing without reconstructing ordinary ordinal multiplication}

For a positive embedded ordinal $\theta$, put
\begin{equation}
 p_\theta(\xi,\zeta)=\theta\xi+\zeta,
 \qquad \Lambda(\theta)=\theta^2+\theta.
 \label{eq:pairing}
\end{equation}
All operations in \eqref{eq:pairing} are field operations, hence natural
ordinal operations on these arguments.

\begin{lemma}[Bounded pairing]
\label{lem:pairing}
The map $p_\theta$ injects $\theta\times\theta$ into $\Lambda(\theta)$.
For every epsilon number $\lambda$ and every $0<\theta<\lambda$,
$\Lambda(\theta)<\lambda$. For infinite $\theta$,
$\card{\Lambda(\theta)}=\card\theta$.
\end{lemma}
\begin{proof}
Suppose $\xi<\xi'<\theta$ and $\zeta,\zeta'<\theta$. Since
$\xi+1\le\xi'$ in ordinal order, positivity and distributivity give
\[
 \theta\xi+\zeta<\theta\xi+\theta
 =\theta(\xi+1)\le\theta\xi'\le\theta\xi'+\zeta'.
\]
For fixed $\xi$, cancellation gives injectivity in $\zeta$. Also
$\theta\xi+\zeta<\theta^2+\theta$. Closure below $\lambda$ follows from its
being a field cutoff. The cardinal estimate is
\cref{lem:card-estimates}, together with $\theta\le\Lambda(\theta)$.
\end{proof}

Consequently a sign sequence of birthday $\Lambda(\theta)$ codes a binary
relation on $\theta$. To make this code canonical for a given relation,
require every plus position to be in the image of $p_\theta$; unused
positions are minus. This last requirement is the first-order formula
\[
 \forall\rho<\Lambda(\theta)\,
 \bigl(\Bit(r,\rho)\longrightarrow
       \exists\xi,\zeta<\theta\ [\rho=p_\theta(\xi,\zeta)]\bigr).
\]
Here and below bounded ordinal quantifiers include the predicate $\Ord$.

For a graph between two different ordinals $\delta,\epsilon<\lambda$, use
$\theta=\delta+\epsilon+1$, which is above both. Coding the graph in the
$\delta\times\epsilon$ rectangle of this square again stays below $\lambda$.
The construction requires no definable global pairing on the entire class
of ordinals, and no reconstruction of ordinary ordinal arithmetic.

\section{Pointed relation codes and the exact interpreted universe}
\label{sec:codes}

\subsection{The definable domain of codes}

A prospective code is a triple
\[
 c=(\delta,r,t),\qquad
 0<\delta<\lambda,\quad t<\delta,\quad\bd(r)=\Lambda(\delta),
\]
with $\delta,t$ embedded ordinals and with unused bits of $r$ equal to minus.
Write
\[
 E_c(i,j)\quad\Longleftrightarrow\quad
 i,j<\delta\ \land\ \Bit(r,p_\delta(i,j)).
\]
The orientation is \emph{child belongs to parent}: $i\,E_c\,j$ will mean
that the object at node $i$ is a member of the object at node $j$.

Require extensionality,
\begin{equation}
 \forall i,j<\delta\,
 \bigl[\forall k<\delta\ (E_c(k,i)\leftrightarrow E_c(k,j))
       \longrightarrow i=j\bigr],
 \label{eq:extensional}
\end{equation}
and well-foundedness in the following form:
\begin{align}
 \forall u\,[\bd(u)=\delta\ \longrightarrow\ (&
 (\exists j<\delta\ \Bit(u,j))\ \longrightarrow\nonumber\\[-0.2em]
 &\exists j<\delta\,[\Bit(u,j)\land
    \forall i<\delta\ (\Bit(u,i)\longrightarrow\neg E_c(i,j))])].
 \label{eq:wf}
\end{align}
By \eqref{eq:fiber-subsets}, this is \emph{actual} well-foundedness of the
set relation, not just well-foundedness against a restricted collection of
subsets.

Finally require generation by the root:
\begin{align}
 \forall u\,[\bd(u)=\delta\ \longrightarrow\ (&
 \Bit(u,t)\land
 \forall i,j<\delta\,[\Bit(u,j)\land E_c(i,j)\longrightarrow\Bit(u,i)]
 \nonumber\\[-0.2em]
 &\longrightarrow\forall i<\delta\ \Bit(u,i))].
 \label{eq:rooted}
\end{align}
Thus no proper downward-closed subset of the graph contains the root.
Let $\Code(c)$ be the conjunction of these conditions. The triple is an
ordinary finite-dimensional interpretation domain in the single-sorted
structure $\B_\lambda$.

\begin{lemma}[Rooted collapse]
\label{lem:collapse}
Every valid code has a unique Mostowski collapse
$\pi_c:\delta\longrightarrow T_c$ onto a transitive set. If
$x_c=\pi_c(t)$, then
\[
 T_c=\TC(\{x_c\}),\qquad
 \card\delta=\card{\TC(\{x_c\})}.
\]
\end{lemma}
\begin{proof}
Define by well-founded recursion
\[
 \pi_c(j)=\{\pi_c(i):i\,E_c\,j\}.
\]
Extensionality and well-founded induction show that $\pi_c$ is injective:
if a collision existed, comparison of predecessor values would give equal
predecessor sets at the first possible collision. Its image is transitive,
and the displayed recursion gives an isomorphism. Uniqueness follows by
well-founded induction as well; this is the usual Mostowski-collapse
argument.

The nodes reachable from $t$ by finitely many downward edges form a
downward-closed subset containing $t$. By \eqref{eq:rooted} this is all of
$\delta$. Under the collapse, these are exactly $x_c$ and its iterated
members. Thus the image is $\TC(\{x_c\})$, giving the cardinal equality.
\end{proof}

Root generation is not needed merely to encode arbitrary sets. It is useful
here because it makes the hereditary-size cost exact rather than just an
upper bound polluted by unrelated graph components.

\subsection{Equality and membership are definable}

For codes $c=(\delta,r,t)$ and $d=(\epsilon,s,u)$, define $c\sim d$ when a
coded bijection $f:\delta\to\epsilon$ preserves $E$, sends $t$ to $u$, and
reflects $E$. Such a graph is expressed by a sign sequence in a square of
width $\delta+\epsilon+1$, as in \cref{lem:pairing}. The familiar function,
bijection, and edge-preservation clauses use only bounded ordinal
quantifiers and bits. Therefore $\sim$ is a single first-order relation.

For membership, let $c\in_\I d$ mean that there is a coded injection
$f:\delta\to\epsilon$ such that
\begin{align}
 &E_d(f(t),u),\label{eq:member-root}\\
 &E_c(i,j)\ \longleftrightarrow\ E_d(f(i),f(j))
       &&(i,j<\delta),\label{eq:member-edge}\\
 &E_d(v,f(j))\ \longrightarrow\
       \exists i<\delta\,[E_c(i,j)\land f(i)=v]
       &&(j<\delta,0\le v<\epsilon).\label{eq:member-full}
\end{align}
In \eqref{eq:member-full}, the quantifier on $v$ simply ranges over the
embedded ordinals below $\epsilon$. The clause says that the image contains
\emph{all} predecessors of its nodes, not merely some of them.

\begin{lemma}[Correctness of equality and membership]
\label{lem:equal-member}
For valid codes,
\[
 c\sim d\iff x_c=x_d,
 \qquad
 c\in_\I d\iff x_c\in x_d.
\]
\end{lemma}
\begin{proof}
A root-preserving isomorphism gives equal collapsed roots by uniqueness of
collapse. Conversely, if the roots agree, the two images are the same
transitive closure by \cref{lem:collapse}, and
$\pi_d^{-1}\pi_c$ gives the required isomorphism. Its graph has size at most
$\max(\card\delta,\card\epsilon)$ and is codable within the same cutoff.

For membership, \eqref{eq:member-edge} and \eqref{eq:member-full}, followed
by induction along $E_c$, imply
$\pi_d(f(j))=\pi_c(j)$ for all $j$. The root condition then gives
$x_c\in x_d$. Conversely, if $x_c\in x_d$, the inclusion
$\TC(\{x_c\})\subseteq\TC(\{x_d\})$ transports through the collapse maps to
an injection satisfying all three clauses. Its image is predecessor-closed,
so \eqref{eq:member-full} really does hold. The graph is again small enough
to be coded below the cutoff.
\end{proof}

Let $\I(\B_\lambda)$ be the quotient of valid codes by $\sim$, with the
membership relation just defined. All quantifiers in its interpreted
formulas are replaced by quantifiers over triples satisfying $\Code$.
No quotient class is assumed to be an element of the original structure.

\subsection{Cardinal rounding}

\begin{theorem}[Exact collapse range]
\label{thm:collapse-range}
For every epsilon number $\lambda$, the map $[c]\mapsto x_c$ is an
isomorphism
\[
 \I(\B_\lambda)\cong(H_{\ceilcard{\lambda}},\in).
\]
More precisely, the ordinals $\delta$ supporting rooted codes for a given
set $x$ are exactly those with
\[
 0<\delta<\lambda,
 \qquad\card\delta=\card{\TC(\{x\})}.
 \label{eq:exact-code-size}
\]
\end{theorem}
\begin{proof}
Every code has $\delta<\lambda\le\kappa(\lambda)$. Hence its collapse has
hereditary size $\card\delta<\kappa(\lambda)$, by
\cref{lem:collapse}. This proves one inclusion of the range.

Conversely, suppose $x\in H_{\kappa(\lambda)}$. Its nonempty transitive
closure has some cardinal $\nu<\kappa(\lambda)$. By the definition of
cardinal ceiling, $\nu<\lambda$: if $\lambda$ is a cardinal this is
immediate; otherwise $\nu\le\card\lambda<\lambda$. Choose a bijection from
$\nu$ onto $\TC(\{x\})$, transfer membership to $\nu$, and designate the
node for $x$. The resulting graph is extensional, well founded, and generated
by that node. It has a sign code of length $\Lambda(\nu)<\lambda$.
Thus $x$ occurs in the range.

For any other $\delta<\lambda$ equipotent to this transitive closure, the
same construction works with a bijection from $\delta$. Necessity in
\eqref{eq:exact-code-size} was already proved by rooted collapse. Correctness
of both relations is \cref{lem:equal-member}.
\end{proof}

\begin{example}[The empty set and a first nonempty code]
The empty set uses one node, no edges, and root zero. With $\delta=1$, the
canonical relation string has birthday $\Lambda(1)=2$ and both signs minus.
Thus a code for the empty set is $(1,--,0)$, not the surreal zero itself.
For $\{\varnothing\}$ take $\delta=2$, root $1$, and the single edge
$0E1$. The relation string has length $6$ and its only plus is at
$p_2(0,1)=1$. Numeric values and the sets they code must not be identified.
\end{example}

\begin{example}[A countable cutoff already recovers all countable sets]
Every hereditarily countable set has a rooted code on a finite ordinal or
on $\omega$. Since
\[
 \Lambda(\omega)=\omega^2+\omega<\varepsilon_0,
\]
all such codes already occur in $\B_{\varepsilon_0}$. They recover
$H_{\omega_1}$, including countable ordinals far larger than
$\varepsilon_0$. Those larger ordinals are recovered as \emph{collapsed
relations}, not as all-plus scalar sign sequences in the cutoff.
\end{example}

\section{Bi-interpretation and recovery of the cutoff}
\label{sec:biinterpretation}

\subsection{The ordinal and sign-graph bridges}

For an actual embedded ordinal $a<\lambda$, there is a distinguished code
$O(a)$ with domain $a+1$, root $a$, and edge relation $iEj$ iff $i<j$.
Its sign relation code is uniquely determined by the unused-bit convention.
The formula defining $O(a)$ is uniform, and its collapse is the
von Neumann ordinal $a$.

\begin{definition}[Visible interpreted ordinals]
For a set code $c$ define
\[
 \Vis(c)\quad\Longleftrightarrow\quad
 \exists a\,[\Ord(a)\land c\sim O(a)].
\]
Under the collapse isomorphism, these are exactly the ordinals below
$\lambda$, not necessarily all ordinals of $H_{\kappa(\lambda)}$.
\end{definition}

There is also a definable bridge $\SG(s,c)$ saying that $c$ represents the
set-coded sign function of the scalar $s$: its domain is $O(\bd(s))$, its
values are the interpreted ordinals $0$ and $1$, and at each actual
$\alpha<\bd(s)$ it has value $1$ exactly when $\Bit(s,\alpha)$ holds.
Application of a set-coded function and the Kuratowski ordered-pair
construction are ordinary fixed formulas in the interpreted membership
language. The ordinal bridge supplies their arguments. Thus $\SG$ is a
fixed first-order formula, not an external assignment operation.

For each $s$ there is exactly one equivalence class of codes $c$ with
$\SG(s,c)$; conversely, every sign function whose length is below $\lambda$
arises this way. Existence follows from the hereditary-size estimate for
sign graphs and \cref{thm:collapse-range}.

\subsection{A first-order test for a cardinal cutoff}

Define the \emph{visible-ordinal collector} formula
\begin{equation}
 \Coll(c)\quad\Longleftrightarrow\quad
 \Code(c)\land
 \forall d\,[\Code(d)\longrightarrow
       (d\in_\I c\leftrightarrow\Vis(d))].
 \label{eq:collector}
\end{equation}
It says that one interpreted set collects all the scalar ordinals. This is
not the assertion that the scalar carrier itself is a set inside itself.

\begin{theorem}[The cutoff is a set exactly in the noncardinal case]
\label{thm:collector}
For an epsilon number $\lambda$,
\[
 \B_\lambda\models\exists c\,\Coll(c)
 \quad\Longleftrightarrow\quad
 \lambda\text{ is not a cardinal}.
\]
When it exists, the collector is unique modulo $\sim$ and its collapse is
exactly the ordinal $\lambda$.
\end{theorem}
\begin{proof}
The elements specified in \eqref{eq:collector} are exactly all ordinals
$a<\lambda$. In the ambient set universe the unique set with precisely
those elements is $\lambda$. Hence a collector exists precisely when
$\lambda\in H_{\kappa(\lambda)}$, equivalently
$\card\lambda<\kappa(\lambda)$. For a cardinal $\lambda$ the right side
fails; otherwise $\kappa(\lambda)=(\card\lambda)^+$ and it holds.
Uniqueness follows from extensionality of the interpreted universe.
\end{proof}

This collector is the additional information needed to recover a
noncardinal cutoff after interpreting its hereditary-size universe. Merely
saying that $\B_\lambda$ interprets $H_{\kappa(\lambda)}$ would lose it.

\subsection{The reverse interpretation}

If $\lambda=\kappa$ is a cardinal, interpret surreals in $H_\kappa$ as all
set-coded functions from an ordinal into $\{0,1\}$. By
\cref{lem:hereditary}, these are exactly the sign sequences of length below
$\kappa$. Define the order by first disagreement, birthday by the all-one
function on the domain ordinal, and arithmetic by the certificate formulas
of \cref{thm:certificates}. Call this interpretation $\J$.

If $\lambda$ is not a cardinal, work instead in the pointed structure
$(H_{\kappa(\lambda)},\in,\lambda)$. Use the same definitions, with the
extra domain restriction that the sign length is below the named ordinal
$\lambda$. The certificate verifier may use auxiliary sign sequences
anywhere in $H_{\kappa(\lambda)}$; it is not required to carry out a large
recursion inside the smaller interpreted sign domain. The root outputs are
below $\lambda$ by classical cutoff closure. This gives a uniform
interpretation $\J_{\rm cut}$ in the language with one named cutoff.

\begin{theorem}[Uniform local bi-interpretation]
\label{thm:biinterpretation}
The interpretations just constructed prove all assertions of
\cref{thm:main-rounding}. The maps witnessing the two bi-interpretations are
definable, not merely external bijections.
\end{theorem}
\begin{proof}
Start on the set side. A code in the interpreted surreal structure gives
an extensional well-founded relation on an ordinal. Its Mostowski-collapse
root defines the map back to the original set. The graph of the collapse
belongs to the same $H_\kappa$: its domain, range, and all values have
hereditary size bounded by a fixed infinite cardinal below $\kappa$.
Consequently this correspondence is defined by a first-order formula in
$H_\kappa$, even when Replacement is not an axiom of that structure.
The collapse equation is a bounded verification of the supplied function.
Every set has a code, and equality and membership are respected by
\cref{lem:equal-member}. In the pointed case, the collector collapses to the
named cutoff, so the correspondence preserves that constant as well.

Start on the surreal side. The formula $\SG(s,c)$ identifies $s$ with the
sign function representing it in the copy of the set universe. In the
cardinal case all sign functions in the interpreted $H_\lambda$ occur.
In the noncardinal case, the collector defines $\lambda$, and the restricted
reverse interpretation uses exactly lengths below that ordinal. Thus $\SG$
is a definable bijection to the twice-interpreted surreal structure.
It preserves numerical order and birthday directly from the sign rules,
and it preserves arithmetic by the certificate theorem. This proves the
required definable isomorphisms in both directions.
\end{proof}

\begin{corollary}[The first epsilon cutoff and hereditary countability]
\label{cor:epsilon0}
The structure $\B_{\varepsilon_0}$ is parameter-free bi-interpretable with
$(H_{\omega_1},\in)$.
\end{corollary}
\begin{proof}
The theorem first gives the pointed structure
$(H_{\omega_1},\in,\varepsilon_0)$. The ordinal $\varepsilon_0$ is
parameter-free definable in $H_{\omega_1}$ as the least nonzero fixed point
of ordinary ordinal $\alpha\mapsto\omega^\alpha$. For any countable
arguments, the functions witnessing ordinal arithmetic and the relevant
well-orders are countable sets, hence belong to $H_{\omega_1}$. Their
well-founded recursive verification is absolute. The least-fixed-point
specification therefore picks out the actual $\varepsilon_0$ there.
Naming it is a definitional expansion, which may be removed.
\end{proof}

The same argument applies to any noncardinal epsilon number that is
parameter-free definable, with its correct ordinal value, in the recovered
hereditary-size universe. It is not an assertion that every ordinal can be
named parameter-free.

\subsection{Coherence of the interpretation}

\begin{proposition}[Coherence under cutoff inclusions]
\label{prop:coherence}
If $\lambda\le\eta$ are epsilon numbers, the inclusion of scalar cutoffs
induces, under $\I$, the actual inclusion
\[
 H_{\kappa(\lambda)}\longrightarrow H_{\kappa(\eta)}.
\]
A code keeps the same collapsed value, and all sign-graph bridges agree.
For cardinal cutoffs the reverse interpretations also commute with the
actual inclusions.
\end{proposition}
\begin{proof}
For a code already below $\lambda$, the larger structure has no additional
sign sequences of birthday $\delta$ for any $\delta<\lambda$: both contain
\emph{all} sequences of each such length in the same ambient universe.
Thus the subset tests in well-foundedness and root generation are unchanged.
Equality and membership agree by their collapse characterizations; any
needed witness graph has a code below $\lambda$ as proved above. The
collapse is defined from the same relation, so its value is unchanged.
The sign-graph bridge refers to the same prefixes and bits. Finally the
reverse arithmetic formulas compute the same external operations.
\end{proof}

If two noncardinal cutoffs have the same cardinal ceiling, the induced
inclusion on recovered \emph{unpointed} set universes is the identity.
Their recovered distinguished cutoffs nevertheless differ. This is precisely
why elementary inclusion will usually fail.

\section{The exact elementary-inclusion spectrum}
\label{sec:elementary}

The elementary-inclusion result is stronger than the observation that a
birthday expansion is not governed by real-closed-field quantifier
elimination. It identifies exactly which initial inclusions are possible.

\begin{lemma}[Noncardinal lower cutoffs never include elementarily]
\label{lem:lower-noncardinal}
If $\lambda<\eta$ are epsilon numbers and $\lambda$ is not a cardinal, then
$\B_\lambda\not\prec\B_\eta$.
\end{lemma}
\begin{proof}
Choose a code $c$ in $\B_\lambda$ for its visible-ordinal collector. Its
collapse is $\lambda$, and it satisfies $\Coll(c)$ there. By coherence it
still collapses to $\lambda$ in $\B_\eta$. But $\lambda$ itself is now a
scalar ordinal, since $\lambda<\eta$. Its interpreted ordinal code is
visible, and is not a member of the set coded by $c$, because
$\lambda\notin\lambda$. Thus $\Coll(c)$ is false in the larger structure.
This single formula, with the three entries of $c$ as parameters, witnesses
failure of elementarity.
\end{proof}

\begin{lemma}[A cardinal lower cutoff cannot include elementarily into a
noncardinal cutoff]
\label{lem:upper-noncardinal}
If $\lambda$ is a cardinal and $\eta$ is not, then
$\B_\lambda\not\equiv\B_\eta$.
\end{lemma}
\begin{proof}
The parameter-free sentence $\exists c\,\Coll(c)$ is false in the first
structure and true in the second, by \cref{thm:collector}.
\end{proof}

\begin{theorem}[Classification of elementary initial inclusions]
\label{thm:elementary}
For epsilon numbers $\lambda<\eta$,
\[
 \B_\lambda\prec\B_\eta
 \quad\Longleftrightarrow\quad
 \lambda,\eta\text{ are cardinals and }H_\lambda\prec H_\eta.
\]
\end{theorem}
\begin{proof}
The preceding two lemmas force both cutoffs to be cardinals. For cardinal
cutoffs, an elementary inclusion of the surreal structures induces an
elementary inclusion of their interpreted set structures. By
\cref{prop:coherence} this is the actual inclusion $H_\lambda\subseteq H_\eta$.
Conversely, an elementary inclusion of those set structures induces an
elementary inclusion of the uniformly interpreted sign structures. The
reverse interpretation and its definable sign-graph isomorphism identify this
with the original inclusion $\B_\lambda\subseteq\B_\eta$.
\end{proof}

\begin{corollary}[Two kinds of elementary extension]
The ordered field inclusion $K_\lambda\subseteq K_\eta$ is elementary for
all epsilon numbers $\lambda<\eta$. Adding birthday changes its elementary
status according to \cref{thm:elementary}. In particular,
\[
 K_{\varepsilon_0}\prec K_{\omega_1},
 \qquad
 \B_{\varepsilon_0}\not\prec\B_{\omega_1}.
\]
\end{corollary}
\begin{proof}
Both fields are real closed, so ordered-field quantifier elimination gives
the first assertion. The birthday assertions are instances of the theorem.
\end{proof}

Despite the last failure, both structures in that example are separately
bi-interpretable with $H_{\omega_1}$. There is no contradiction: a
bi-interpretation does not make \emph{every particular inclusion} elementary.
The smaller structure defines an interpreted cutoff set; that formula changes
its meaning under the displayed inclusion.

\begin{corollary}[Comparison with the full class]
For an epsilon number $\lambda$, interpreted formula by formula,
\[
 \B_\lambda\prec\B_{\mathrm{all}}
 \quad\Longleftrightarrow\quad
 \lambda\text{ is a cardinal and }H_\lambda\prec V.
\]
\end{corollary}
\begin{proof}
The same collector obstruction excludes a noncardinal lower cutoff. For a
cardinal lower cutoff, the uniform interpretation and reverse sign-graph
interpretation extend to definable classes. The induced set map is
$H_\lambda\subseteq V$, so the elementary equivalence follows for each
formula. This uses no internal satisfaction class for $V$.
\end{proof}

The condition $H_\lambda\prec H_\eta$ in the theorem is not replaced here by
a conjectural large-cardinal characterization. It is an exact reduction of
one elementary-inclusion problem to another familiar one, together with the
complete exclusion of every noncardinal initial cutoff.

\section{Birthday growth and set-existence axioms}
\label{sec:growth}

\subsection{Birthday eternity}

\begin{definition}[The eternity scheme $\Etern$]
For every formula $\varphi(a,x,\vec p)$ in the birthday-field language,
include the assertion that if $\gamma$ is an embedded ordinal and
\[
 \forall a<\gamma\ \exists!x\ \varphi(a,x,\vec p),
\]
then some embedded ordinal $\beta$ satisfies
\[
 \forall a<\gamma\ \forall x\,
       (\varphi(a,x,\vec p)\longrightarrow\bd(x)<\beta).
\]
The tuple $\vec p$ consists of arbitrary parameters from the structure.
\end{definition}

This is a precisely specified version of birthday boundedness for definable
ordinal-indexed functions. The analogy with the eternity principle in the
Chen--Hamkins--Yang announcements is intentional; the cutoff classification
below is proved directly, rather than imported from an unexpanded axiom
scheme in the slides.

\begin{theorem}[Eternity forces cardinality of an epsilon cutoff]
\label{thm:eternity-noncardinal}
If $\lambda$ is a noncardinal epsilon number, then
$\B_\lambda\not\models\Etern$.
\end{theorem}
\begin{proof}
Put $\mu=\card\lambda$, so $\mu<\lambda$ and
$\kappa(\lambda)=\mu^+$. Choose a bijection $f:\mu\to\lambda$.
Its graph has hereditary size at most $\mu$, so it belongs to
$H_{\mu^+}$ and has a code in $\B_\lambda$. Fix one such code as a
parameter. Using the interpreted graph of $f$ and the ordinal bridge, a
single formula now defines a scalar map on the actual ordinal $\mu$:
at $a<\mu$ its value is the embedded ordinal $f(a)<\lambda$.
Its birthdays exhaust $\lambda$, hence have no bound below $\lambda$.
This violates the corresponding instance of $\Etern$.
\end{proof}

At a singular \emph{cardinal} $\kappa$, an external cofinal map into
$\kappa$ generally has hereditary size $\kappa$ and need not belong to
$H_\kappa$. Thus the same argument cannot simply be repeated with
$\mu=\operatorname{cf}(\kappa)$. Whether an internally definable cofinal
map exists is the point of the next theorem.

\subsection{The axioms that all hereditary-size universes already satisfy}

\begin{lemma}[The fixed background axioms]
\label{lem:H-background}
For every uncountable cardinal $\kappa$, $H_\kappa$ satisfies
Extensionality, Empty Set, Pairing, Union, Infinity, Foundation, every
Separation instance, and Choice. Moreover the actual transitive closure of
each of its elements belongs to $H_\kappa$, and its cardinality is correctly
computed there.
\end{lemma}
\begin{proof}
Transitivity gives Extensionality and Foundation. Empty Set, Pairing, Union,
and Infinity preserve hereditary size below $\kappa$. For any fixed formula
and parameters, the externally defined subset of a set $a\in H_\kappa$
satisfying that formula in the set structure $H_\kappa$ has hereditary size
bounded by that of $a$ (up to finitely many additional objects). It therefore
belongs to $H_\kappa$, proving Separation, including formulas with unbounded
quantifiers in that structure.

An ambient choice function for a family $a\in H_\kappa$ has its domain and
chosen values inside the transitive closure of $a$, and its graph has the
same infinite hereditary-cardinality bound. It belongs to $H_\kappa$.
This proves Choice for set families, not a class-choice principle.

Writing $T=\TC(\{x\})$, the transitive closure of $T$ adds at most the set
$T$ to its already transitive data, so $T\in H_\kappa$. Its description as
the least transitive set containing $x$ as an element is correct internally:
the actual least such set is available there. A bijection between $T$ and
its least equipotent ordinal has hereditary size below $\kappa$, and any
bijection asserted internally is externally a bijection. Minimality among
ordinals is therefore also computed correctly.
\end{proof}

\subsection{Replacement is exactly eternity at cardinal cutoffs}

\begin{theorem}[Eternity--Replacement equivalence]
\label{thm:eternity-replacement}
For every uncountable cardinal $\kappa$,
\[
 \B_\kappa\models\Etern
 \quad\Longleftrightarrow\quad
 H_\kappa\models\text{every instance of Replacement}.
\]
\end{theorem}
\begin{proof}
Assume first Replacement in $H_\kappa$. A definable scalar function on an
ordinal $\gamma<\kappa$ corresponds under the bi-interpretation to a
set-theoretically definable function assigning to $a<\gamma$ its unique
sign graph. Replacement collects those sign graphs in one set of
$H_\kappa$. The union of their domains, with one added successor for each
length, is bounded below $\kappa$: otherwise the transitive closure of the
collected set would contain at least $\kappa$ ordinals. More explicitly,
\[
 \beta=\sup\{\dom(s)+1:s\text{ is an image sign graph}\}<\kappa
\]
bounds their birthdays strictly. This proves the corresponding eternity
instance; the empty-domain case is vacuous.

Conversely, assume eternity and consider a formula defining, in
$H_\kappa$, a function $F:a\to H_\kappa$, where $a\in H_\kappa$ and the
formula has fixed parameters. Choose a bijection $e:\delta\to a$ with
$\delta<\kappa$. Its graph belongs to $H_\kappa$. For each $i<\delta$, put
externally
\[
 \nu_i=\card{\TC(\{F(e(i))\})},
\]
regarded as an initial ordinal. By \cref{lem:H-background}, this is a
uniquely defined ordinal of $H_\kappa$, with the correct external value.
The definition of $F$, the parameter $e$, and the hereditary-cardinality
formula translate into a definable scalar map $i\mapsto\nu_i$ in
$\B_\kappa$. Here the ordinal bridge is available for every ordinal below
$\kappa$. No definable choice of a code for $F(e(i))$ is required: its
\emph{hereditary cardinal} is unique.

Eternity gives a $\beta<\kappa$ bounding all these ordinals. Ambient
Replacement forms the actual image set $A=F[a]$, and its transitive closure
has cardinality at most
\[
 \max(\aleph_0,\card\delta,\card\beta)<\kappa.
\]
Indeed it is contained, up to the image set itself, in the union of
$\card\delta$ transitive closures, each of cardinality at most
$\max(\aleph_0,\card\beta)$. Thus $A\in H_\kappa$, proving the desired
Replacement instance in that structure. This proof uses a fixed cardinal
bound below $\kappa$; it never assumes that an arbitrary union of fewer than
$\kappa$ small sets is small.
\end{proof}

\begin{corollary}
If $\kappa$ is uncountable and regular, then $\B_\kappa\models\Etern$.
\end{corollary}
\begin{proof}
Any externally given function on an ordinal below $\kappa$ has fewer than
$\kappa$ birthdays in its image. Regularity bounds them below $\kappa$.
This argument applies in particular to definable functions.
\end{proof}

\subsection{Packing a birthday fiber into one code}

\begin{definition}[The fiber-packing principle $\Pack$]
For every embedded ordinal $\alpha$, there exist a positive embedded ordinal
$\theta$ and a scalar $r$, with
\[
 \chi=\theta+\alpha+1,
 \qquad\bd(r)=\Lambda(\chi),
\]
such that
\begin{equation}
 \forall u\,[\bd(u)=\alpha\longrightarrow
   \exists\xi<\theta\ \forall\zeta<\alpha\,
    (\Bit(u,\zeta)\leftrightarrow\Bit(r,p_\chi(\xi,\zeta)))].
 \label{eq:packing}
\end{equation}
\end{definition}

This is one first-order sentence, using the previously defined bit and
pairing formulas. Its rows list every sign string of birthday $\alpha$;
repetitions are allowed. Unused coordinates can be ignored. It does not use
ordinary ordinal concatenation or a primitive power-set operation.

\begin{theorem}[Exact packing threshold]
\label{thm:packing}
In $\B_\lambda$, the instance of \eqref{eq:packing} for a particular
$\alpha<\lambda$ holds exactly when
\[
 2^{\card\alpha}<\kappa(\lambda),
 \quad\text{equivalently, }2^{\card\alpha}<\lambda.
\]
where the cardinal on the left is compared as an ordinal in the second
inequality. Consequently $\Pack$ holds exactly when $\lambda$ is an
uncountable strong-limit cardinal.
\end{theorem}
\begin{proof}
A packing table with $\theta<\lambda$ rows gives a surjection from $\theta$
onto $\Pow(\alpha)$, so $2^{\card\alpha}\le\card\theta<\kappa(\lambda)$.
Conversely, if that power-set cardinal is $\nu<\kappa(\lambda)$, then
$\nu<\lambda$ by cardinal rounding. Choose an enumeration of the fiber by
$\theta=\nu$ and encode its table in the square of width
$\chi=\theta+\alpha+1$. The square code stays below $\lambda$ by closure,
so the required $r$ exists. The equivalence of the two inequalities follows
because no cardinal lies in $[\lambda,\kappa(\lambda))$.

For a cardinal $\lambda$, validity for all $\alpha<\lambda$ is exactly the
strong-limit property. If $\lambda$ is not a cardinal, let
$\mu=\card\lambda<\lambda$ and use $\alpha=\mu$. Cantor's theorem gives
$2^\mu\ge\mu^+=\kappa(\lambda)$, so this instance fails.
\end{proof}

\begin{proposition}[Power Set in the interpreted universe]
\label{prop:powerset}
For an uncountable cardinal $\kappa$,
\[
 H_\kappa\models\text{Power Set}
 \quad\Longleftrightarrow\quad
 \kappa\text{ is strong limit}.
\]
For an individual $\alpha<\kappa$, the actual power set $\Pow(\alpha)$ is
an element of $H_\kappa$ exactly when $2^{\card\alpha}<\kappa$.
\end{proposition}
\begin{proof}
Every subset of $\alpha$ is already in $H_\kappa$, so an internal power set
of $\alpha$, if it exists, is its full actual power set. Its hereditary
size is $2^{\card\alpha}$ up to harmless finite or countable bounds.
This proves necessity and the individual assertion. Conversely, for
$x\in H_\kappa$, the transitive closure of $\Pow(x)$ has size at most
\[
 \max(\aleph_0,\card{\TC(\{x\})},2^{\card x}).
\]
This is below a strong-limit $\kappa$, so every required power set exists
in $H_\kappa$.
\end{proof}

Notice that for a noncardinal cutoff $\lambda$, the recovered
$H_{\kappa(\lambda)}$ and the scalar domain do not have the same ordinal
height. Assertions about \emph{all} interpreted ordinals and assertions
about only scalar ordinals must be kept separate. The failure at
$\alpha=\card\lambda$ is enough to settle the full packing principle.

\section{Worldly cutoffs and separation examples}
\label{sec:worldly}

\subsection{The exact conjunction}

Call an uncountable cardinal $\kappa$ \emph{worldly} here when
$V_\kappa\models\mathrm{ZFC}$. This definition does not impose regularity.
We need the following familiar set-theoretic fact, for which a proof is
included to fix the precise meaning.

\begin{lemma}[When hereditary size becomes cumulative rank]
\label{lem:worldly-HV}
For an uncountable cardinal $\kappa$,
\[
 H_\kappa\models\mathrm{ZFC}
 \quad\Longleftrightarrow\quad
 H_\kappa=V_\kappa\text{ and }V_\kappa\models\mathrm{ZFC}.
\]
Moreover, $V_\kappa\models\mathrm{ZFC}$ by itself implies
$H_\kappa=V_\kappa$.
\end{lemma}
\begin{proof}
Suppose $H_\kappa\models\mathrm{ZFC}$. Its internal cumulative hierarchy
exists to every ordinal $\alpha<\kappa$. Every subset of each of its sets
belongs to $H_\kappa$, so its internal power sets are actual power sets.
Induction on $\alpha$ shows that its internal $V_\alpha$ is the actual
$V_\alpha$; the limit step is ordinary union over the actual smaller
ordinals. Thus every $V_\alpha$ for $\alpha<\kappa$ belongs to $H_\kappa$,
and transitivity gives $V_\kappa\subseteq H_\kappa$.
The reverse inclusion is \cref{lem:hereditary}.

Conversely, assume $V_\kappa\models\mathrm{ZFC}$. For every
$\alpha<\kappa$, that model can well-order $V_\alpha$ and supply a bijection
with some ordinal $\delta<\kappa$. These are actual sets and actual
bijections, because the model is transitive. Hence
$\card{V_\alpha}<\kappa$. The transitive closure of any
$x\in V_\kappa$ is contained in some such $V_\alpha$, and therefore has
size below $\kappa$. Again $H_\kappa=V_\kappa$ follows. The remaining
implication is immediate from equality.
\end{proof}

\begin{theorem}[Worldly-cutoff characterization]
\label{thm:worldly}
For an epsilon number $\lambda$,
\begin{align*}
 \B_\lambda\models\Etern+\Pack
 &\iff \lambda\text{ is a cardinal and }H_\lambda\models\mathrm{ZFC}\\
 &\iff \lambda\text{ is a worldly cardinal}.
\end{align*}
For regular uncountable $\lambda$, this holds exactly when $\lambda$ is
strongly inaccessible.
\end{theorem}
\begin{proof}
Either growth principle excludes noncardinal epsilon cutoffs. At cardinal
cutoffs, \cref{thm:eternity-replacement,prop:powerset,lem:H-background}
identify the two missing ZFC requirements as Replacement and Power Set.
The equivalence with the rank formulation is \cref{lem:worldly-HV}.
For regular uncountable cardinals, eternity is automatic and packing is
exactly strong limit, giving the usual inaccessible condition.
\end{proof}

This proves \cref{thm:main-growth}. It does not identify our two formulas with
every formulation of a theory called SA. It is an exact theorem about the
explicit schemes and sentence defined in this article.

\subsection{The two growth principles do different work}

\begin{example}[Eternity without fiber packing]
At $\kappa=\omega_1$, regularity gives eternity. The birthday-$\omega$ fiber
has $2^{\aleph_0}\ge\aleph_1$ members, so no ordinal below $\omega_1$ can
index a packing of it. Hence
\[
 \B_{\omega_1}\models\Etern,
 \qquad\B_{\omega_1}\not\models\Pack.
\]
In the decoded structure, every individual real is present, but the set of
all reals is not an element of $H_{\omega_1}$.
\end{example}

\begin{example}[Fiber packing without eternity]
\label{ex:bethomega}
Let $\kappa=\beth_\omega=\sup_{n<\omega}\beth_n$, with
$\beth_0=\aleph_0$ and $\beth_{n+1}=2^{\beth_n}$. This is a strong-limit
cardinal: for every $\mu<\kappa$ some $\beth_n$ is at least $\mu$, and then
$2^\mu\le\beth_{n+1}<\kappa$. Thus packing holds.

In $H_\kappa$, each finite initial segment of the beth recursion is present
and correctly computed. The necessary finite sequence of power sets,
cardinals, and bijections has hereditary size below $\kappa$, and Power Set
is correct there by \cref{prop:powerset}. A single formula therefore defines
$n\mapsto\beth_n$ on the actual $\omega$. Transporting this formula through
the ordinal bridge gives a definable scalar function whose birthdays are
cofinal in $\kappa$. Consequently
\[
 \B_{\beth_\omega}\models\Pack,
 \qquad\B_{\beth_\omega}\not\models\Etern.
\]
The infinite sequence of all these ordinals is precisely what cannot be
collected as an element of $H_\kappa$.
\end{example}

\begin{example}[A noncardinal field cutoff]
Both principles fail at $\varepsilon_0$. Packing fails at the ordinal
$\omega$. Eternity fails by \cref{thm:eternity-noncardinal}; one may also use
the definable finite power-tower sequence cofinal in $\varepsilon_0$,
with ordinary ordinal exponentiation computed in the decoded
$H_{\omega_1}$. Nevertheless the field is real closed and its birthday
expansion is bi-interpretable with hereditary countability.
\end{example}

These are actual structures exhibiting independence of the two displayed
growth requirements within this family. They are not an assertion of
independence from an unspecified list of all the other axioms of SA.

\subsection{Singular cutoffs satisfying both principles}

\begin{proposition}[A conditional source of singular worldly cutoffs]
\label{prop:singular-worldly}
If $\Omega$ is strongly inaccessible, there is a cardinal $\kappa<\Omega$
of cofinality $\omega$ such that
\[
 V_\kappa\prec V_\Omega.
\]
In particular, $\kappa$ is worldly and $\B_\kappa$ satisfies both growth
principles.
\end{proposition}
\begin{proof}
Enumerate all existential formulas in the language of set theory. Construct
strictly increasing infinite cardinals $\alpha_n<\Omega$. Given
$\alpha_n$, consider the first $n$ formulas and all their finite parameter
tuples from $V_{\alpha_n}$. Whenever $V_\Omega$ has a witness, choose one.
There are fewer than $\Omega$ such requirements, since
$\card{V_{\alpha_n}}<\Omega$. Regularity bounds their witness ranks below
$\Omega$. Choose a cardinal $\alpha_{n+1}<\Omega$ above that bound and above
$\alpha_n$.

Set $\kappa=\sup_n\alpha_n$. Regularity of $\Omega$ gives $\kappa<\Omega$;
it is a cardinal of cofinality $\omega$. Every finite tuple in $V_\kappa$
appears in one stage, and every existential formula is treated at all
sufficiently late stages. Thus every existential statement true in
$V_\Omega$ with parameters in $V_\kappa$ has a witness in $V_\kappa$.
The Tarski--Vaught criterion gives elementarity. Since $V_\Omega$ satisfies
ZFC, so does $V_\kappa$. Apply \cref{thm:worldly}.
\end{proof}

The reflection construction is standard set-theoretic reasoning, included
as an illustration rather than as a novelty claim. It demonstrates why
replacing ``worldly'' by ``inaccessible'' in \cref{thm:worldly} would be wrong.
No inaccessible cardinal is assumed in the unconditional cutoff theorems.

\section{Outer models, new subsets, and a sharp parameter bound}
\label{sec:extensions}

\subsection{Coherence when the sign fibers themselves grow}

Assume $M\subseteq N$ are transitive models of ZFC with the same ordinals
(or the same ordinal height for set models). Let $\kappa$ be an uncountable
cardinal of both models. The structures $\B_\kappa^M$ and $\B_\kappa^N$
use their respective internal sign sequences. Unlike two cutoffs in a single
universe, they may have different birthday fibers even at the same small
ordinal.

\begin{proposition}[Cross-model coherence]
\label{prop:cross-model}
There are actual inclusions
\[
 \B_\kappa^M\subseteq\B_\kappa^N,
 \qquad H_\kappa^M\subseteq H_\kappa^N,
\]
compatible with the interpretations and their definable isomorphisms.
Consequently
\[
 \B_\kappa^M\prec\B_\kappa^N
 \quad\Longleftrightarrow\quad H_\kappa^M\prec H_\kappa^N.
\]
\end{proposition}
\begin{proof}
A sign sequence of $M$ remains the same sign sequence in $N$; order and
birthday are absolute. A small arithmetic certificate existing in $M$
remains valid in $N$, so the field operations agree. An old set of hereditary
size below $\kappa$ comes with an old bijection witnessing that bound; since
$\kappa$ is still a cardinal, it belongs to $H_\kappa^N$.

For code validity, one must not claim that the two models have the same
subsets of the coding ordinal. Instead, a relation well founded in $M$ has
an ordinal rank function in $M$, which remains decreasing in $N$. Its
well-foundedness therefore persists despite the extra subsets in $N$.
Root generation says that every node is connected to the root by a finite
downward path, as established in \cref{lem:collapse}; these paths are
absolute. Extensionality is bounded and absolute. Failure of any of the
code conditions in $M$ also persists, using its old witness. Thus validity
on old codes is unchanged.

The old Mostowski collapse still satisfies its bounded defining equations
in $N$ and hence is the same collapse. Equality and membership of decoded
old sets are unchanged. Witnesses to their equality and membership can be
constructed in $M$ from those collapses, so their required graph codes
already exist there. The reverse sign-graph interpretation likewise sends
old objects to their actual images. Applying each interpretation to an
elementary inclusion now proves the equivalence.
\end{proof}

\subsection{The old-power-set detector}

\begin{theorem}[Uniform detector for an added subset]
\label{thm:detector}
Suppose $\alpha<\kappa$, $N$ has a subset of $\alpha$ not in $M$, and
\[
 \nu=(2^{\card\alpha})^M<\kappa.
\]
Then $\B_\kappa^M\not\prec\B_\kappa^N$. One fixed interpreted formula
witnesses the failure using parameters coding $\alpha$ and
$A=\Pow^M(\alpha)$.
\end{theorem}
\begin{proof}
The cardinal bound gives $A\in H_\kappa^M$. In the membership language
consider
\begin{equation}
 \exists u\,[\forall v\in u\ (v\in\alpha)\ \land\ u\notin A].
 \label{eq:new-subset-detector}
\end{equation}
This is false in $H_\kappa^M$, because all its subsets of $\alpha$ lie in
$A$. It is true in $H_\kappa^N$: any newly added subset of $\alpha$ has
hereditary size below the still-cardinal $\kappa$ and is not an element of
the old set $A$. Translate \eqref{eq:new-subset-detector} through $\I$ and
choose old codes for the two parameters. Cross-model coherence gives the
claimed failure in the birthday structures.
\end{proof}

The membership-language formula \eqref{eq:new-subset-detector} is
$\Sigma_1$: its matrix has only bounded quantifiers. Its fully expanded
birthday-field translation need \emph{not} be existential or $\Sigma_1$ in
that different language, because code validity includes the bit-quantified
well-foundedness test. We make no quantifier-complexity claim across that
translation.

\begin{proposition}[Sharpness of the detector's parameter cost]
\label{prop:detector-cost}
The old parameter $\Pow^M(\alpha)$ belongs to $H_\kappa^M$ exactly when
$(2^{\card\alpha})^M<\kappa$. In our rooted coding scheme, every code for
that parameter has node-domain cardinality
\[
 \card{\TC(\{\Pow^M(\alpha)\})}^M,
\]
which equals the power-set cardinal for infinite $\alpha$. In that infinite case, a code of
birthday cardinality no larger than that cardinal can be chosen.
\end{proposition}
\begin{proof}
The first statement is \cref{prop:powerset} inside $M$. The exact
node-domain cardinality is \eqref{eq:exact-code-size}. For infinite
$\alpha$, the transitive closure consists of the power set, its subsets,
and the ordinal data beneath $\alpha$, giving cardinality
$2^{\card\alpha}$. Choosing its initial ordinal as the node domain gives
relation birthday $\Lambda(\nu)$, of cardinality $\nu$ by
\cref{lem:pairing}.
\end{proof}

This is not a theorem that elementarity must hold when the bound fails.
Another formula can detect a change at a smaller cutoff. It says exactly
when this particular, universally applicable old-power-set witness is
available.

\begin{example}[A conditional continuum-sized illustration]
Suppose $M$ satisfies CH, $N$ preserves $\omega_2^M$, and $N$ adds a real.
Then the old set of reals has hereditary size $\omega_1^M$, so the detector
works at the cutoff $\omega_2^M$. That old set is not an element of
$H_{\omega_1}^M$, so the same parameter cannot be used at the cutoff
$\omega_1^M$. No claim of elementarity at the smaller cutoff follows.
\end{example}

\subsection{Rigidity under extensions at a strong limit}

\begin{theorem}[Strong-limit extension rigidity]
\label{thm:extension-rigidity}
Under the hypotheses of \cref{prop:cross-model}, assume in addition that
$\kappa$ is strong limit in $M$. Then
\begin{align*}
 \B_\kappa^M\prec\B_\kappa^N
 &\iff \B_\kappa^M=\B_\kappa^N\\
 &\iff \Pow^M(\alpha)=\Pow^N(\alpha)
                 \quad\text{for every }\alpha<\kappa.
\end{align*}
\end{theorem}
\begin{proof}
If some subset of an $\alpha<\kappa$ is new, the strong-limit property
makes its old power-set parameter small enough, and \cref{thm:detector}
excludes elementarity. Conversely, if all these power sets agree, every
sign sequence in $N$ of length $\alpha<\kappa$ has its plus-set in $M$.
Reconstructing the sign sequence from that plus-set in $M$ shows that the
two scalar domains agree. Their operations agree by
\cref{prop:cross-model}, so the structures are equal and hence the
inclusion is elementary. Equality of the scalar domains also forces
equality of the corresponding plus-sets, proving the remaining direction.
\end{proof}

This proves \cref{thm:main-extension}. It applies, for example, to any
forcing extension satisfying the stated cardinal-preservation hypotheses,
but uses no particular forcing construction. It is a theorem about
\emph{standard cutoff structures at a fixed common cardinal}, not a
prohibition on abstract elementary extensions supplied by model theory.

Their ordered field reducts still include elementarily even when the
birthday expansion fails the detector, since both reducts are real closed.
Thus a set-theoretic change can be invisible to ordered-field formulas yet
visible to one fixed birthday-field formula with old parameters.

\section{Surcomplex numbers and the orientation obstruction}
\label{sec:surcomplex}

\subsection{The natural enriched surcomplex structure}

For an epsilon number $\lambda$, put $C_\lambda=K_\lambda[i]$. This is the
usual pair field, with conjugation $\sigma(x+iy)=x-iy$. Introduce the
coordinate-birthday function
\[
 \beta(x+iy)=\max(\bd(x),\bd(y)),
\]
whose value is an embedded real ordinal. Consider the two structures
\[
 \C_\lambda^-=(C_\lambda;0,1,+,-,\cdot,\sigma,\beta),
 \qquad
 \C_\lambda^+=(\C_\lambda^-,i).
\]
Here a choice of square root $i$ of $-1$ is named in the second structure
but not in the first. These are specific enrichments of the surcomplex
field, not claims about the field language alone.

\begin{proposition}[Recovery of the real birthday structure]
\label{prop:real-recovery}
Both surcomplex structures interpret $\B_\lambda$ without parameters.
The structure $\C_\lambda^+$ is uniformly bi-interpretable with
$\B_\lambda$.
\end{proposition}
\begin{proof}
The real carrier is the definable fixed field of $\sigma$. Its order is
recovered by
\[
 x<y\quad\Longleftrightarrow\quad
 \exists t\,[\sigma(t)=t\land t\ne0\land t^2=y-x],
\]
using real closedness. The restriction of $\beta$ is exactly $\bd$.
This gives the first interpretation.

Conversely, interpret complex numbers as pairs of reals, with
\[
 (x,y)(u,v)=(xu-yv,xv+yu),\qquad
 \sigma(x,y)=(x,-y),\qquad i=(0,1),
\]
and with $\beta$ the maximum of the two birthdays. In the structure with
$i$ named, the map $(x,y)\mapsto x+iy$ is definable, with inverse
\[
 z\longmapsto
 \left(\frac{z+\sigma(z)}2,
       \frac{z-\sigma(z)}{2i}\right).
\]
These give the two definable isomorphisms required for a bi-interpretation.
\end{proof}

Thus cardinal rounding, recovery of a noncardinal cutoff, and the resulting
set-theoretic distinctions transfer to the oriented surcomplex structures.
In particular the interpretations do not acquire a new set-theoretic
size requirement merely by adjoining $i$.

\subsection{Why omitting the imaginary unit changes bi-interpretability}

\begin{lemma}[Rigidity of a standard birthday cutoff]
\label{lem:rigidity}
Every automorphism of $\B_\lambda$ is the identity.
\end{lemma}
\begin{proof}
The embedded ordinals are definable and numerically well ordered. Any
order automorphism of this ordinal segment fixes every ordinal, by
transfinite induction. The automorphism therefore fixes all birthday
values. The bit predicate is definable, so it preserves the sign at each
fixed ordinal coordinate. Every sign sequence is determined by its domain
and bits, and hence is fixed.
\end{proof}

This is a standard simplicity/birthday rigidity argument, not a new
classification of the unrestricted field automorphisms studied elsewhere
in the repository \cite{RepoAutomorphisms}.

\begin{theorem}[The two-element orientation obstruction]
\label{thm:orientation}
The automorphism group of $\C_\lambda^-$ consists exactly of the identity
and conjugation. Consequently $\C_\lambda^-$ is mutually interpretable
with, but is not parameter-free bi-interpretable with, the corresponding
set structure
\[
 (H_{\kappa(\lambda)},\in)
 \quad\text{or, in the noncardinal case,}\quad
 (H_{\kappa(\lambda)},\in,\lambda).
\]
After naming $i$, the bi-interpretation holds as in
\cref{prop:real-recovery,thm:biinterpretation}.
\end{theorem}
\begin{proof}
An automorphism preserves the fixed field of $\sigma$, its definable order,
and the birthday function on it. Its restriction is therefore the identity
by \cref{lem:rigidity}. The image of $i$ is a root of $X^2+1$, so it is
$i$ or $-i$. Since every element has the form $x+iy$, the automorphism is
respectively the identity or conjugation. Both do preserve $\sigma$ and
$\beta$.

Each transitive membership structure $H_\kappa$ has no nontrivial
automorphism: if all members of $x$ are fixed, extensionality fixes $x$,
and external well-founded induction applies to every $x$. Naming an ordinal
does not change this rigidity. Were $\C_\lambda^-$ bi-interpretable with
such a structure, its nontrivial conjugation would induce the identity on
the interpreted set structure, hence the identity on the twice-interpreted
surcomplex structure. The definable isomorphism back to
$\C_\lambda^-$ would then force conjugation itself to be the identity,
a contradiction.

Mutual interpretation still holds: recover the real birthday structure
inside $\C_\lambda^-$ and then its set universe; in the opposite direction
use the reverse surreal interpretation and the pair field, simply omitting
the distinguished constant $i$. The missing definable choice of orientation
is exactly the obstruction to upgrading this to a bi-interpretation.
\end{proof}

The group computation is a consequence of classical rigidity. The point for
this article is the precise distinction between mutual interpretability and
bi-interpretability, and why an unnamed imaginary unit cannot be treated as
harmless notation in the latter claim.

\section{What the theorems do and do not imply}
\label{sec:scope}

\subsection{Logical strength versus field strength}

Already $\B_{\varepsilon_0}$ recovers $H_{\omega_1}$. Within that membership
structure, the natural numbers and all subsets of the natural numbers form
the usual full two-sorted structure of second-order arithmetic: each subset
is hereditarily countable, and arithmetic on finite ordinals is definable.
Thus the birthday-field interpretation includes quantification over all
actual subsets of $\omega$, not merely over computable subsets. Its ordinary
first-order language can carry that quantification because the numbers
coding those subsets are first-order elements.

This does not make set theory definable in the pure ordered-field reduct.
The needed birthday structure is absent there. Nor does it give a decision
procedure: a hypothetical effective enumeration of all true birthday-field
sentences would, by the effective interpretation of standard arithmetic,
enumerate all true arithmetic sentences, which is impossible. These are
logical-complexity consequences, not new algorithms for deciding truth.

\subsection{Four statements that must not be conflated}

For a cutoff $\lambda$ the following can have very different answers:
\begin{center}
\begin{tabular}{@{}P{0.31\textwidth}P{0.61\textwidth}@{}}
\toprule
Question & Exact issue in this article\\
\midrule
Is $K_\lambda$ a real closed field? & Classical epsilon-number cutoff theorem.\\[0.35em]
Which sets can be decoded? & Exactly $H_{\kappa(\lambda)}$, by rooted relation codes.\\[0.35em]
Does the decoded universe satisfy ZFC? & At cardinal cutoffs, exactly the worldly condition.\\[0.35em]
Is an initial inclusion elementary? & Both cutoffs must be cardinals, followed by elementarity of their hereditary-size universes.\\
\bottomrule
\end{tabular}
\end{center}
An algebraic closure theorem alone does not answer any of the last three
questions. Conversely, a transitive set model satisfying ZFC need not arise
from a regular inaccessible cutoff.

\subsection{Choice and quotients}

Ordinary Choice is used to enumerate an individual transitive closure and,
in some proofs, to enumerate a particular fiber or choose a particular
bijection. The interpretation does not choose one representative
simultaneously for every set. It keeps all valid codes and quotients by a
definable equivalence relation. This is why Global Choice is unnecessary.

The exact cardinal-rounding conclusion is stated in ZFC. Without Choice,
not every transitive closure is well-orderable, and a code on an ordinal
cannot automatically represent every set in a proposed hereditary-size
universe. Removing Choice therefore requires changing the target or the
coding language. No choice-free version is claimed here.

\subsection{Set models, nonstandard models, and class truth}

At a fixed uncountable cardinal, the structures in the main theorems are
actual set structures, with ordinary first-order satisfaction. Their
interpretations are not hypothetical syntactic models satisfying only a
fragment of surreal arithmetic. A nonstandard model of ZFC can internally
perform the same constructions, but its coded well-foundedness and its
internal cardinal ceiling must be read in that model. In particular, an
externally ill-founded model is not covered by an external Mostowski-collapse
claim merely because it internally believes its code is well founded.

For the full class $\No$, one applies the translation to a fixed formula
and fixed set parameters. Statements about ``all formulas'' are metatheoretic
schemes. We do not form a set of all class functions, a proper-class
membership quotient, or a truth predicate for $V$ inside ZFC. The cutoff
results themselves do not require a class theory, a Grothendieck universe,
or a large-cardinal existence axiom.

\subsection{Why uncountability and the epsilon hypothesis matter}

For the main hereditary-size bi-interpretation theorem at cardinal cutoffs,
uncountability ensures that countably sized computation certificates and the
ordinary infinite set $\omega$ remain below the cutoff. At $\omega$, the
finite-birthday surreals form the dyadic ring rather than a field; it would
be incorrect to reuse the real-closed-field argument there.

For noncardinal cutoffs, the epsilon hypothesis supplies closure of the
bounded pairing terms and field operations. This article does not classify
arbitrary ordinal sign cutoffs that fail those closure conditions. One can
formulate relation-code interpretations under weaker closure hypotheses,
but the full birthday-field statement would then need a different language
or partial operations.

\section{Research audit and further mathematical targets}
\label{sec:audit}

\subsection{Dependency and novelty ledger}

\begin{longtable}{@{}P{0.25\textwidth}P{0.37\textwidth}P{0.30\textwidth}@{}}
\toprule
Result & Main inputs & Status here\\
\midrule\endhead
Global birthday-field interpretation & Pointed well-founded relation coding & Prior Chen--Hamkins--Yang announcement; not claimed new.\\[0.5em]
Real-closed epsilon cutoffs & Classical surreal field theory & Prior van den Dries--Ehrlich, with correction.\\[0.5em]
Small arithmetic certificates & Prefix option graphs, finite words, graph rank & Explicit local proof ingredient; no historical priority claim for the general locality principle.\\[0.5em]
Cardinal rounding and cutoff recovery & Bounded pairing, rooted collapse, sign-graph bridge & Proposed local refinement, including noncardinal and singular cases.\\[0.5em]
Elementary-inclusion spectrum & Collector formula and coherent bi-interpretations & Proposed exact classification for initial cutoff inclusions.\\[0.5em]
Eternity and packing spectrum & Hereditary cardinal bounds; internal correctness of small cardinalities & Proposed quantitative calibration; growth/set-existence analogy is prior work.\\[0.5em]
Old-power-set detector & A fixed membership-language formula and code-cost equality & Proposed explicit refinement; ordinary outer-model set reasoning is standard.\\[0.5em]
Orientation obstruction & Real-axis interpretation and classical rigidity & A consequence clarifying the surcomplex bi-interpretation statement, not a new unrestricted automorphism classification.\\
\bottomrule
\end{longtable}

The central place where the argument could not be replaced by a general
slogan is \cref{sec:certificates}. Saying that the global construction is
set-theoretically definable would not prove that the same formula computes
arithmetic in $H_\kappa$ when that structure lacks Replacement. The
certificate bound resolves that issue. A second essential point is the
collector: a noncardinal cutoff is not recovered merely by interpreting the
unpointed hereditary-size universe.

\subsection{Scope of the source inspection}

The public sources inspected for the current global interpretation program
were Hamkins's September and November 2025 announcements and slides and the
March 2026 seminar announcement. The available slides explicitly mark the
work as new and its details as being checked. This article treats the
announced result as prior work while supplying its own local proofs; it does
not claim to have audited an unpublished complete manuscript.

The repository comparison used the pinned report catalog, the foundations
report's relevant construction, recursion, universe, and cutoff discussions,
and the surcomplex-automorphism report's stated scope. Repository search did
not provide a dependable exhaustive full-text index, so a negative keyword
search was not used as proof of absence. No complete build or formal theorem
inventory of the repository was performed. No repository files were changed.

Literature searches specifically combined surreal arithmetic with bounded
birthdays, hereditary-size universes, cardinal cutoffs, bi-interpretation,
and worldly cardinals. No inspected source stated the exact local spectrum
proved here. This supports a \emph{proposed novelty} assessment, not a
categorical claim that no specialist has previously obtained the results.

\subsection{Questions now reduced to more precise form}

One target is to replace the remaining condition $H_\kappa\prec H_\tau$
in \cref{thm:elementary} by useful explicit hypotheses in interesting
families of cardinals. The theorem already rules out every noncardinal
initial inclusion, so the residual question is genuinely about hereditary
set universes, not field embeddings.

A second target is a syntactic comparison with a finalized axiom system for
surreal arithmetic. The principles $\Etern$ and $\Pack$ are deliberately
explicit here. Establishing exact definitional equivalence with particular
published axiom formulations should be a separate theorem, rather than an
identification by terminology.

A third target is formalization of the size-certified interpretation. The
necessary objects are ordinary set-coded sign functions, bounded relation
codes, well-founded option graphs, and quotient interpretations. The first
milestone should be the pairing and bit lemmas; the main proof obligation is
the certificate verifier and its hereditary-size bound. The present article
contains no Lean proof terms and claims no kernel verification.

A fourth target is the choice-free analogue. Ordinal-indexed well-founded
relations naturally recover sets with well-orderable transitive closures.
An exact account in ZF would have to state that restriction and examine
whether a richer surreal coding signature can escape it. Nothing in the
ZFC proofs justifies silently deleting Choice.

\section{Conclusion}

Birthday truncation has a precise set-theoretic effect. At an epsilon-number
cutoff $\lambda$, rooted sign codes recover the hereditary-size universe at
the cardinal ceiling of $\lambda$. When $\lambda$ is not itself a cardinal,
its value survives as a definable interpreted set. This extra datum explains
why noncardinal cutoff inclusions never become elementary, even though all
their ordered field reducts are elementary subfields.

The growth calculations distinguish two independent resources. Packing a
birthday fiber requires its full power-set cardinal to fit below the cutoff.
Bounding definable ordinal-indexed birthday growth corresponds, at cardinal
cutoffs, to collecting definable images. Their conjunction detects worldly
cardinals, not merely inaccessible ones. Under an outer-model extension,
the old power set supplies an equally explicit detector of added subsets,
with an exact hereditary-size cost.

The same information is available in surcomplex arithmetic once conjugation,
coordinate birthday, and a choice of imaginary unit are specified. Omitting
the last choice preserves mutual interpretation but leaves a genuine
two-element automorphism obstruction to bi-interpretation. Across these
results, the decisive structure is not the field alone: it is the field
with a carefully specified memory of how its elements are born.

\appendix
\section{Formula expansion and size checks}
\label{app:formulas}

\subsection{What ``first-order'' means in the coding proof}

Every abbreviation used above expands to a finite formula in
$\{0,1,+,-,\cdot,<,\bd\}$. Bounded ordinal quantifiers are shorthand for
ordinary scalar quantifiers with guards $\Ord(x)$ and $x<\delta$.
A set variable in an interpreted membership formula becomes a triple of
scalar variables. Equality becomes $\sim$, and membership becomes
$\in_\I$; quantifiers are relativized to $\Code$.

For example, a graph code $f$ on $\delta\times\epsilon$, embedded in the
square of width $\theta=\delta+\epsilon+1$, uses
\[
 F(i,j)\iff i<\delta\land j<\epsilon\land
                 \Bit(f,p_\theta(i,j)).
\]
It is a total function exactly when
\[
 \forall i<\delta\ \exists j<\epsilon\,
 \bigl(F(i,j)\land\forall k<\epsilon\ (F(i,k)\longrightarrow k=j)\bigr).
\]
Injectivity is
\[
 \forall i,i'<\delta\ \forall j<\epsilon\,
       (F(i,j)\land F(i',j)\longrightarrow i=i'),
\]
and surjectivity is $\forall j<\epsilon\,\exists i<\delta\,F(i,j)$.
Occurrences of $f(i)$ in the main text can therefore be eliminated by these
unique-value clauses. Root preservation and edge correspondence add only
finitely many quantifiers.

The graph code's birthday is $\Lambda(\theta)$; all plus positions outside
the relevant rectangle can be required absent. The same convention makes
canonical relation coding a definable operation when the intended relation
is already specified by a formula and its bit string exists.

\subsection{Every potentially large construction has a stated bound}

\begin{longtable}{@{}P{0.35\textwidth}P{0.56\textwidth}@{}}
\toprule
Object & Bound used in the proof\\
\midrule\endhead
Subsets tested in well-foundedness & Individual sign strings of birthday $\delta$; no collection of that whole fiber is required.\\[0.4em]
A relation on a code domain $\delta$ & Sign birthday $\Lambda(\delta)=\delta^2+\delta<\lambda$.\\[0.4em]
An isomorphism or member embedding & Square width $\delta+\epsilon+1<\lambda$, with the corresponding $\Lambda$ bound.\\[0.4em]
Collapse value of a rooted code & Transitive closure has exactly $\card\delta$ members.\\[0.4em]
An arithmetic computation & At most $\mu$ word vertices and hereditary size $\mu$, where $\mu$ bounds the inputs and is infinite.\\[0.4em]
Rank labels in that computation & Below $\mu^+$, hence each has cardinality at most $\mu$.\\[0.4em]
Replacement image after eternity & One bound $\max(\aleph_0,\card\delta,\card\beta)<\kappa$, not an arbitrary singular-cardinal union.\\[0.4em]
A packed birthday-$\alpha$ fiber & Exists exactly if $2^{\card\alpha}<\kappa(\lambda)$.\\[0.4em]
The old-power-set parameter & Hereditary size $(2^{\card\alpha})^M$ for infinite $\alpha$.\\
\bottomrule
\end{longtable}

\subsection{An essential negative regression}

In the membership definition, edge preservation alone is insufficient.
Take the code for $x=\varnothing$ and the three-node chain coding
$y=\{\{\varnothing\}\}$. Map the sole node of the first code to the middle
node of the chain, which codes $\{\varnothing\}$. The image root is a child
of the target root, and the one-node induced subgraph has no edges, just
like the source. But $\varnothing\notin\{\{\varnothing\}\}$.
The predecessor-full clause \eqref{eq:member-full} rejects this false
witness, because the middle node has a predecessor omitted from the image.

\section{Finite verification and reproduction}
\label{app:verification}

The accompanying script \path{finite_checks.py} checks finite analogues of
the prefix formula, the bounded pairing map, rooted relation coding,
isomorphism, and the membership test. It also checks the negative regression
in the preceding appendix. Its concrete counts and outcome are recorded in
\path{verification_report.json}. The script uses only the Python standard
library and may be run by
\begin{quote}\ttfamily python3 finite\_checks.py\end{quote}

The delivered run passed 16,129 ordered sign-pair tests on all 127 strings
of length at most six, 642 sign-coordinate tests, and 11,440 finite pairing
checks. It examined all 530 directed relations on domains of size one
through three, all 1,570 choices of a rooted relation, and all 225 ordered
pairs among the 15 valid rooted extensional well-founded codes, for both
equality and membership. The deliberately incomplete membership test
accepted the false witness above; the predecessor-full version rejected it.

These checks can expose a reversed edge, a missing root condition, a bad
pairing index, or an incorrect finite sign formula. They do not prove
transfinite well-foundedness, cardinal bounds, bi-interpretability,
Replacement equivalences, or any large-cardinal assertion. The proofs of
those statements are the written arguments above.

The PDF is built from this standalone source; the bibliography is included
below and no external graphics or font files are needed. Run
\begin{quote}\ttfamily
pdflatex -interaction=nonstopmode -halt-on-error surreal\_cutoffs.tex
\end{quote}
three times, or use \texttt{latexmk -pdf}. Document-production checks and
finite tests are separate from formal mathematical verification.

\begin{thebibliography}{99}

\bibitem{Conway}
J. H. Conway, \emph{On Numbers and Games}, second edition,
A K Peters, 2001. Classical number forms and recursive arithmetic.

\bibitem{Gonshor}
H. Gonshor, \emph{An Introduction to the Theory of Surreal Numbers},
London Mathematical Society Lecture Note Series 110,
Cambridge University Press, 1986. Sign sequences and the simplicity theorem.

\bibitem{vdDE}
L. van den Dries and P. Ehrlich, ``Fields of surreal numbers and
exponentiation,'' \emph{Fundamenta Mathematicae} \textbf{167} (2001),
173--188. Read with the following erratum.

\bibitem{vdDEerr}
L. van den Dries and P. Ehrlich, ``Erratum to `Fields of surreal numbers and
exponentiation','' \emph{Fundamenta Mathematicae} \textbf{168} (2001),
295--297.

\bibitem{BG}
O. Bournez and Q. Guilmant, ``Surreal fields stable under exponential and
logarithmic functions,'' arXiv:2201.08199, 2022.
\url{https://arxiv.org/abs/2201.08199}.
The corrected cutoff and real-closedness statements used here are stated
in Theorems 1.1--1.2; their original source is \cite{vdDE,vdDEerr}.

\bibitem{RM}
D. R. Rangel and H. L. Mariano, ``An algebraic (set) theory of surreal
numbers, I,'' arXiv:1911.12726, 2019.
\url{https://arxiv.org/abs/1911.12726}.
Foundational and algebraic set-theory background, distinct from the
birthday-field interpretation considered here.

\bibitem{CHYKobe}
J. D. Hamkins, reporting joint work with J. Chen and R. Yang,
``The elementary theory of surreal arithmetic is bi-interpretable with set
theory,'' Kobe, September 2025; announcement dated August 20, 2025,
with linked slides.
\url{https://jdh.hamkins.org/theory-of-surreal-arithmetic-bi-interpretable-with-set-theory-japan-september-2025/}.
The consulted slides announce the global result and discuss the coding and
growth principles; they are not a complete refereed paper.

\bibitem{CHYNotreDame}
J. D. Hamkins, reporting joint work with J. Chen and R. Yang,
``The elementary theory of surreal arithmetic is bi-interpretable with set
theory,'' Notre Dame Logic Seminar, November 2025, with linked slides.
\url{https://jdh.hamkins.org/surreal-arithmetic-notre-dame-logic-seminar-nov-2025/}.

\bibitem{CHYCUNY}
New York Logic Workshop, ``Surreal arithmetic is bi-interpretable with set
theory,'' seminar announcement for J. D. Hamkins, March 13, 2026,
reporting joint work with J. Chen and R. Yang.
\url{https://nylogic.github.io/logic-workshop/2026/03/13/surreal-arithmetic-is-bi-interpretable-with-set-theory.html}.

\bibitem{RepoCatalog}
V. Reshetnikov, \emph{Surreal} repository, documentation catalog,
\path{docs/README.md}, inspected at commit
\texttt{4f2645599121fa872c7104995e47f86f0382351f} on September 22, 2026.
\url{https://github.com/VladimirReshetnikov/Surreal/tree/4f2645599121fa872c7104995e47f86f0382351f/docs}.
Repository research manuscripts are treated as prior project material,
not as independently refereed publications.

\bibitem{RepoFoundations}
\emph{Foundations for Surreal and Surcomplex Mathematics: Sets, Classes,
Universes, Type Theories, and a Pinned Lean Audit}, maintained repository
exposition, September 2026. At the commit in \cite{RepoCatalog}, see
\path{docs/foundations-and-computation/foundations/article.tex}
and its README. Relevant inspected topics include sign presentations,
local recursion, universes, and birthday cutoffs.

\bibitem{RepoAutomorphisms}
\emph{Automorphisms of the Surcomplex Numbers: Real Forms, Hahn Symmetries,
Valuation, and Rigidity}, repository research exposition,
September 22, 2026. At the commit in \cite{RepoCatalog}, see
\path{docs/surcomplex/surcomplex-field-automorphisms/README.md}.
The scope statement was inspected to avoid treating broad automorphism
rigidity as a new contribution of the present article.

\end{thebibliography}
\end{document}
